\subsection{3.3. 上下文范式}\label{the-context-paradigm}

第 3.1 节与第 3.2 节各自作用于一个上下文：前者作为效应的载体，后者作为余效应的载体，而同时承载两者的单一上下文应是什么样的，则尚未明确。本节为这一统一给出具体构造，从余效应组装出一个观测等价性，以提供第 3.1.3 节所遗留的效应独立性，并论证所得的上下文类型本身就构成一种编程范式。

\subsection{3.3.1. 统一上下文}\label{unified-context}

对于上下文 $\Gamma$，效应上下文 $\partial \Gamma$（第 3.1 节）提供了一种更高层的抽象：它携带上一层的上下文以及该层级的累加器（定义~2）。把这一结构递归化，并与余效应上下文 $\Sigma$ 相结合，得到如下类型：

定义 32. 上下文类型 \(\Gamma _ { \infty }\) 定义如下：

\[
\Gamma _ { \infty } : = \mu \Gamma . \Gamma \times ( \Gamma \to \Gamma ) \times \Sigma\tag{31}
\]

其中三个投影分别为：

• \(\Gamma\)：当前上下文状态（递归的）；

• \(\Gamma \to \Gamma\)：累加器，用于恢复该层级的效应；

• \(\Sigma\)：携带依赖信息的余效应上下文。

在该定义下，效应把 \(\mathfrak { E } _ { \Gamma _ { \infty } }\) 映射到自身，从而把 $\partial$-塔统一为单一的自相似类型。余效应上下文 $\Sigma$ 在结构上被整合进来：依赖操作（set、get）作用于 $\Sigma$，而累加器追踪它们的逆转。由于支撑 $\Sigma$ 的类型族 $\mathcal{V}$ 不受约束，系统需要在组件间共享的任何状态，都可以编码为带有适当值类型的依赖——$\Sigma$ 涵盖所有共享可变状态，而不只是组件间的依赖。组件与其环境之间的每次交互，都穿过这一单一实体。

层次化组合。\(\Gamma _ { \infty }\) 的递归结构支持层次化控制：父上下文聚合多个子层级的效应，形成树状的控制结构，在保持模块化的同时支持统一的跨层级管理。效应变换实现了一种字面意义上的“即插即用”隐喻：

• 加载组件对应于执行其效应（插入）；

• 卸载组件对应于恢复其效应（拔出，且不影响其他运行中的组件）；

• 层级中不同层级的组件可独立加载与卸载；父上下文聚合并管理其所有子级的效应，从而实现任意深度的嵌套组合。

\subsection{3.3.2. 观测等价性}\label{observational-equivalence}

第 3.1 节的恢复保证断言状态之间的一种相等（定理~7），这是一种理想化，因为物理状态无法恢复到其原状。例如，free 向分配器释放一个块，却不会恢复堆在 malloc 之前所具有的布局；生成式名字也不会被丢弃它的逆变换所恢复，因为下一次创建会抽取一个全新的名字 [39]。因此，第 3 节的相等关系都应在某个等价关系 ≃ 之下理解，而我们取 ≃ 为观测等价性：两个状态当且仅当没有任何观测者能区分它们时才相关。比较行为而非表示，是建立程序等价的既定路径 [40]，而这样的比较所得出的关系，取决于观测者所能利用的信息 [41]。上下文观测者所能利用的，是它所携带的余效应；每个余效应都自带各自的等价关系（定义~24），因此上下文上的关系由这些余效应的等价关系组装而成。组装这一关系正是本小节的工作，而按它作商，正是获得第 3.1.3 节所要求的独立性之所在。

定义 33. 两个余效应上下文当它们把相同的键绑定到相关的值上时相关；两个上下文状态当它们的余效应投影相关时相关：

\[
\begin{array} { l l l } { { \sigma \simeq \sigma ^ { \prime } } } & { { : = } } & { { \mathrm { d o m } ( \sigma ) = \mathrm { d o m } ( \sigma ^ { \prime } ) \wedge \forall k \in \mathrm { d o m } ( \sigma ) . \sigma ( k ) \simeq _ { k } \sigma ^ { \prime } ( k ) } } \\ { { \gamma \simeq \gamma ^ { \prime } } } & { { : = } } & { { \sigma _ { \gamma } \simeq \sigma _ { \gamma ^ { \prime } } } } \end{array}\tag{32}
\]

其中 \(\sigma _ { \gamma }\) 表示 \(\gamma\) 的余效应投影（定义~32）。

没有被任何键绑定的那部分状态因此被遗忘，而正是这种遗忘使得定理~7 能在 ≃ 之下被阅读：上述例子中的堆布局与生成式名字都落在这个关系之外，除非有某个键绑定它们。第 3.2.2 节对 ≃ 所需要的性质是随之而来的，而非被假定的。相关状态具有相同的定义域，因此在满足谓词 \(\sigma \models d\) 以及定义 26 的分类判定 \(\mathrm { notify } _ { d }\) 上一致，且响应性是 \(\Sigma / \simeq\) 的一个性质。

把这一关系称为观测性的，是对每个 \(\simeq _ { k }\) 的一种主张，即它所区分的，不会多于键 \(k\) 的操作所能区分的。值的观测者运行这些操作并读取它们的输出。

定义 34. 设 \(V\) 按定义~24 的意义携带一个操作集合 \(\mathcal{A}\)，并把效应函数 \(a ( x )\) 在所有参数 \(x : X _ { a }\) 上的变换幺半群（定义~17）记为 \(\mathfrak { M } ( a )\)。\(\mathcal{A}\) 上的一个测试，是幺半群 \(\mathfrak { M } ( a )\)（\(a \in \mathcal { A }\)）的生成元之上的有限字，其中每个字母施加于它前面各字母所留下的值之上；它的输出是那些为操作的前向映射的字母沿途产生的输出；在某个前提条件不满足处，它未定义。当 \(\mathcal{A}\) 上的每个测试要么在 \(v\) 与 \(v ^ { \prime }\) 两处都已定义、要么两处都未定义，且在两处产生相同的输出时，值 \(v , v ^ { \prime } : V\) 不可区分，记作 \(v \approx _ { \mathcal { A } } v ^ { \prime }\)。

引理 35. 不可区分性是操作所尊重的最粗关系。即：

\begin{enumerate}
\def\labelenumi{\arabic{enumi}.}
\item
  按定义~24 的意义，\(\mathcal{A}\) 的每个操作都尊重 \(\approx _ { \mathcal { A } }\)；
\item
  \(\mathcal{A}\) 的每个操作所尊重的每个等价关系，都包含于 \(\approx _ { \mathcal { A } }\)。
\end{enumerate}

因此，\(\simeq _ { k }\) 的每个可容许选择都包含于 \(\approx _ { \mathcal { A } _ { k } }\)，而 \(\approx _ { \mathcal { A } _ { k } }\) 本身是可容许的。

证明.

\begin{enumerate}
\def\labelenumi{\arabic{enumi}.}
\item
  设 \(v \approx _ { \mathcal { A } } v ^ { \prime }\)，并设 \(a \in \mathcal { A }\) 施加于某个参数。在一个测试前加一个字母仍然是测试，因此前向映射所到达的值不可区分，任一已产生的逆从不可区分的参数所到达的值也不可区分；单字母测试给出两处都已定义或两处都未定义，且输出相等。
\item
  设 \(R\) 是这样一个等价关系，且 \(v R v ^ { \prime }\)。测试的每个字母都是某个操作的前向映射或已产生的逆，而尊重关系沿任一路径携带 \(R\)，使所到达的值保持相关、每个字母处的输出保持相等。因此每个测试在 \(v\) 与 \(v ^ { \prime }\) 处结果一致。 □
\end{enumerate}

处处用 ≃ 替换 = 本身并不足够，因为效应函数除了状态之外还返回逆变换，而 ≃ 所等同的两个状态，必须产生 ≃ 也等同的逆变换。

定义 36. 当满足以下条件时，映射 \(f : \Gamma \to \Gamma\) 尊重 ≃：

\[
\begin{array} { r } { \forall \gamma , \gamma ^ { \prime } \in \Gamma . \quad \gamma \simeq \gamma ^ { \prime } \Rightarrow f ( \gamma ) \simeq f ( \gamma ^ { \prime } ) } \end{array}\tag{33}
\]

两个映射当它们在每个状态处一致时相关；\(\partial \Gamma\) 中的两个对当两个分量都相关时相关：

\[
\begin{array} { r c l } { f \simeq g } & { : = } & { \forall \gamma \in \Gamma . \ f ( \gamma ) \simeq g ( \gamma ) } \\ { ( \delta , g ) \simeq ( \delta ^ { \prime } , g ^ { \prime } ) } & { : = } & { \delta \simeq \delta ^ { \prime } \wedge g \simeq g ^ { \prime } } \end{array}\tag{34}
\]

尊重 ≃ 的映射是能下降到 \(\Gamma / \simeq\) 的映射，由 ≃ 相关的两个映射是在那里下降到同一个映射的两个映射。效应函数两者都需要：前者使它所计算的状态在商上被确定，后者使它返回的逆变换在商上被确定。

定义 37. 在 ≃ 之下重读定义~8：当 \(e\) 作为映射 \(\Gamma \to \partial \Gamma\) 尊重 ≃，并且（记 \(( \delta , g ) = e ( \gamma )\)）对每个 \(\gamma \in \Gamma\) 满足

\begin{enumerate}
\def\labelenumi{\arabic{enumi}.}
\tightlist
\item
  \(g ( \delta ) \simeq \gamma ;\)
\item
  \(g\) 尊重 ≃。
\end{enumerate}

时，\(e \in \mathfrak { E } _ { \Gamma }\) 属于 \(\mathfrak { E } _ { \Gamma } ^ { * }\)。取 ≃ 为 \(\Gamma\) 上的相等关系，即恢复定义~8。

引理 38. 把 \(\mathfrak { E } _ { \Gamma } ^ { * }\) 按定义~37 理解，第 3.1 节所断言的每个状态相等，在把 = 换成 \(\simeq\) 之后依然成立，而且从 \(( \gamma _ { 0 } , \mathrm { i d } _ { \Gamma } )\) 出发可达的每个状态的累加器都尊重 \(\simeq\)。

证明. 累加器是各逆变换的复合，其中每个逆变换按定义~37(2) 尊重 ≃，而尊重 ≃ 的映射的复合也尊重 \(\simeq\)，基例是 \(\mathrm { i d } _ { \Gamma }\)。第 3.1 节的证明于是原样成立——尊重关系正是使关系穿过逆变换的东西：从 \(g _ { 2 } ( \delta _ { 2 } ) \simeq \delta _ { 1 }\) 与 \(g _ { 1 } ( \delta _ { 1 } ) \simeq \gamma\) 出发，尊重给出 \(( g _ { 1 } \circ g _ { 2 } ) ( \delta _ { 2 } ) \simeq \gamma\)，这正是每次复合逆变换所采取的步骤；而定理~7 的可靠性不变式据此步骤读作 \(\varphi ( \gamma ) \simeq \gamma _ { 0 }\)。 □

定义~19 所要求的交换性，由同一引理在 ≃ 之下理解，而正是这样理解，才使交换性变得可以达成：两个操作可能留下 \(\simeq _ { k }\) 所等同的值，却仍算作可交换。对两个操作，它比其提升所诱导的效应函数多要求一件事：操作还会产生输出。

定义 39. 当操作 \(a\) 与 \(a ^ { \prime }\) 的提升在每对参数处都作为效应函数（定义~19）相互独立，且任何一方的变换都不扰乱另一方所产生的输出时，它们独立：

\[
\forall x : X _ { a } , g \in \mathfrak { M } ( a ^ { \prime \Sigma } ) , \sigma \in \Sigma . \quad \operatorname { p r } _ { 3 } \big ( a ^ { \Sigma } ( x ) ( g ( \sigma ) ) \big ) = \operatorname { p r } _ { 3 } \big ( a ^ { \Sigma } ( x ) ( \sigma ) \big )\tag{35}
\]

交换 \(a\) 与 \(a ^ { \prime }\) 后同样成立，其中按定义~34 用 \(\mathfrak { M } ( a ^ { \Sigma } )\) 表示提升 \(a ^ { \Sigma } ( x )\) 在所有参数上的变换幺半群，正如 \(\mathfrak { M } ( a )\) 表示操作本身的变换幺半群。当 \(\mathcal { A } _ { k }\) 中任意两个操作都独立（一个操作也视为与自身独立）时，键 \(k\) 是可交换的。

对于不同键，该条件直接成立。

定理 40. 位于不同键上的操作是独立的。

证明. 设 \(a\) 位于 \(\mathcal { A } _ { k }\)，\(a ^ { \prime }\) 位于 \(\mathcal { A } _ { k ^ { \prime } }\)，且 \(k \neq k ^ { \prime }\)。按定义~24，\(\mathfrak { M } ( a ^ { \Sigma } )\) 的每个生成元都具有 \(\sigma \mapsto \sigma [ k \mapsto u ( \sigma ( k ) ) ]\) 的形式，其中 \(u\) 是 \(\mathcal { V } _ { k }\) 上的映射，它或是前向映射的提升，或是已产生的逆的提升；\(a ^ { \prime }\) 在 \(k ^ { \prime }\) 处同理。这样的两个映射可交换——每个都只读写一个键，且两个键互不相同——而引理~18(1) 把交换性从生成元扩展到两个幺半群。对于第二个条件，\(a ^ { \Sigma }\) 在 \(\sigma\) 处所产生的——逆变换与输出都一样——由 \(\sigma ( k )\) 决定，而 \(\mathfrak { M } ( a ^ { \prime \Sigma } )\) 的每个生成元都保持它原样。 □

值是一张可独立增删条目的表的键是可交换的，路由或事件监听器的注册就是代表性情形：以任意顺序进行两次注册，都会留下一张对每个测试都给出相同回答的表，而任一注册都可以在另一注册仍在场时被撤销。值是一个有序链的键则不可交换，因为插入在另一个中间件之前的中间件会看到不同的请求，而且两种顺序中任一种都无法在不扰乱另一种的情况下被撤销。开篇例子的分配器依其接口所发布的内容而分野：当它所发出的句柄不被该键的任何操作所比较时，\(\simeq _ { k }\) 可以在句柄重命名之下关联两个堆——这正是 CompCert 关联程序与其翻译的记忆状态的方式 [42]——此时分配是可交换的；当地址是依相等性比较的输出时，没有任何可容许的 \(\simeq _ { k }\) 能使两种分配顺序一致，此时该键不可交换。

组件所执行的是一系列操作，其中每个操作都可能依赖其前面操作所产生的输出，而下面这条定理所说的，正是这种形状的效应函数。

定义 41. 由余效应中介的效应函数构成最小的集合 \(\mathfrak { E } _ { \Sigma } ^ { \mathcal { A } } \subseteq \mathfrak { E } _ { \Sigma }\)，它包含单位 \(\eta _ { \Sigma }\)，并在下述构造下封闭：对于键 \(k\)、操作 \(a \in \mathcal { A } _ { k }\)、参数 \(x : X _ { a }\) 以及成员族 \(( e _ { b } ) _ { b \in B _ { a } }\)，

\[
\sigma \mapsto \mathbf { l e t } ~ ( \delta , s , b ) = a ^ { \Sigma } ( x ) ( \sigma ) \ \mathbf { i n } \ \mathbf { l e t } ~ ( \varepsilon , t ) = e _ { b } ( \delta ) \ \mathbf { i n } ~ ( \varepsilon , s \circ t )\tag{36}
\]

仍然是一个成员。每个阶段执行一个操作，并按输出选择其后继，因此某个参数可以依赖已经获得的输出。一个成员中出现的操作，就是其各阶段在每种输出选择下所执行的操作。

定理 42. 设 \(e _ { 1 } , e _ { 2 } \in \mathfrak { E } _ { \Sigma } ^ { \mathcal { A } }\)，并设两者操作均出现的每个键都是可交换的（定义~39）。则 \(e _ { 1 }\) 与 \(e _ { 2 }\) 独立（定义~19）。

证明. 对定义~41 的构造作归纳，\(\mathfrak { M } ( e _ { i } )\) 位于由 \(e _ { i }\) 中出现的操作的生成元所生成的子幺半群中：单位生成平凡幺半群，而一个阶段是 \(a ^ { \Sigma } ( x )\) 与某个成员的 ⋄-复合，对此可应用引理~18(2)。

对定义~19 的条件 (1)，由引理~18(1)，只需 \(e _ { 1 }\) 中出现的某个操作的生成元与 \(e _ { 2 }\) 中出现的某个操作的生成元可交换即可。当两个操作位于不同键时，这就是定理~40；当它们位于同一个键时，该键同时承载两者的操作，按假设是可交换的。

对条件 (2)，取 \(g \in \mathfrak { M } ( e _ { 2 } )\)——它是 \(e _ { 2 }\) 中出现的操作的生成元的复合——并对 \(e _ { 1 }\) 的构造作归纳。单位在每个状态处给出 \(\mathrm { i d } _ { \Sigma }\)。在一个阶段处，设 \(( \delta , s , b ) = a ^ { \Sigma } ( x ) ( \sigma )\)，\(( \varepsilon , t ) = e _ { b } ( \delta )\)，则该阶段在 \(\sigma\) 处给出 \(s \circ t\)。操作之间的独立性——每次只施加于 \(g\) 的一个生成元——在 \(g ( \sigma )\) 处再次给出 \(s\) 与 \(b\)，因此选择同一个后续计算 \(e _ { b }\)，而条件 (1) 使它所运行的状态位于 \(g ( \delta )\) 处，归纳假设在那里再次给出 \(t\)。该阶段因此在 \(g ( \sigma )\) 处给出 \(s \circ t\)。 □

组件与其环境的每次交互都穿过上下文，而类型族 \(\mathcal{V}\) 不受约束，因此系统可以把跨组件共享的每个位置都绑定到它自己的一个键上（第 3.3.1 节）。组件的效应函数于是是由余效应中介的效应函数沿余效应投影的提升，独立性转移到该提升上，而它的变换只移动投影本身。第 3.1.3 节遗留的假设由此得到满足，随之而来的，是整个组件系统的时序可组合性。

该分解所划分的，是计算的交换部分与其顺序敏感部分。交换部分由效应承载：组件按其任务所需的任意顺序执行它们，推论~21 按系统认为方便的任意顺序逆转它们，任意两个组件都不互相约束。顺序敏感部分由余效应承载：因为操作不可交换的键，其顺序必须从效应外部施加，而施加顺序有两个位置可用。在单一组件内部，累加器施加顺序，无论效应为何都按 LIFO 顺序逆转（定理~16）。跨组件则由声明的余效应施加顺序：一个组件提供另一个组件所声明的余效应，提供先于声明的满足（第 3.2.2 节）。可组合性由此在组件而非单个效应的粒度上获得，这正是第 4 节所工作的尺度。

该定理有两个限度值得指明。把每个共享位置都绑定到一个键，是这一范式的纪律，而非该构造的性质；因此，系统无法具体化为余效应的位置，落在第 6.1 节的边界之外，也随之落在定理之外。而键的可交换性是键所发布的接口的性质，因此满足它是提供该键的组件的义务，而非消费它的组件的义务。

\subsection{3.3.3. 上下文范式的定位}\label{situating-the-context-paradigm}

编程范式在处理副作用的方式上有着根本差异。两个公认的极点界定了这一谱系：

显式状态线程化（函数式）。为保持引用透明性，纯函数式语言把副作用建模为状态上的显式变换。State monad \(S \to ( A , S )\) [23] 使环境穿过每一次计算。这一方法给出很强的组合保证：效应在类型中可见，且适合等式推理。然而，它带来显著的可用性代价：调用链中的每个函数都必须接受并返回状态参数，即使它只是原样传递状态。随着效应维度的增长（日志、配置、I/O），monad 堆叠或 effect-handler 样板代码不断膨胀。

隐式修改（命令式/面向对象）。主流命令式语言允许组件修改共享状态并访问依赖，而无须在调用点显式声明。在效应一侧，代表性例子是 React 的 useEffect hook：它在组件内部的 fiber 上注册一个持久的副作用，但效应目标与注册机制都不作为显式参数出现——识别依赖的是隐藏运行时状态中的调用顺序位置。在余效应一侧，Java 的服务定位器模式（例如 Spring 的 ApplicationContext.getBean(\ldots)）在运行时从进程级注册表获取依赖，在每个调用点都需要判空与类型转换；依赖关系是隐式的，散落于整个代码库。更一般地，理解 f() 如何修改或依赖系统，需要传递性地阅读其实现。重构因此变得脆弱，因为移动或删除一次调用，可能悄然破坏远处的不变式。

上下文范式结合了函数式方法的可追踪性与命令式方法的可用性。效应与余效应都通过一个显式的上下文参数来中介。因此，每个操作都可归因于它被调用于其上的那个具体上下文，进而归因于该上下文所属的组件。

除结合两个极点的长处之外，上下文范式还让开发者逐个处理每个效应与依赖，并把它们自动组合进系统的行为之中。对于可逆转的效应，开发者只需提供每个原子操作的逆变换，任何复合的逆变换则按组合得出（第 3.1 节），因此组件的拆卸是由其装载推导而来，而非与之并列书写。对于响应式的余效应，组件只需声明它需要的依赖，运行时自动解析并重新连线它们（第 3.2 节），在提供者被添加、移除或替换时保持连线一致。在两条路径上，原本依赖于开发者纪律的正确性，都成为范式的结构性性质。
